\begin{lemma}[Outside-swap normalization of an exit triple]
Every \(\noderef{TchangeCounterexample}\) \(C\) can be replaced, without
changing its \(\noderef{TchangeCounterexampleMeasure}\), by a
counterexample \(D\) whose three legs have
been relabeled so that \(b_1\prec_D b_3\) and \(b_2\prec_D b_3\), and
for which
\[
 D.\varphi=\noderef{OutsideTreeColorSwap}
   (C.T,C.\varphi,C.a,D.a).
\]
Here \(D.T=C.T\), \(D.b=C.b\), and the new first color \(D.a\) is missing
at \(D.b_3\) under the new coloring \(D.\varphi\).
\end{lemma}
\begin{proof}
The three tree ends of \(C\) are distinct.  Their displayed indices in
the tree-sequence are therefore distinct, so the linear order encoded by
\(\noderef{TreeVertexOrder}\) has a unique largest one.  Relabel the three
legs so that this largest end is called \(b_3\); call the other two
\(b_1,b_2\).  This permutation preserves every counterexample condition
and the sum of the three lengths.

Choose a color \(\gamma\) missing at \(b_3\).  Such a color exists because
the \(\noderef{KTriple}\) condition gives
\(k\geq\Delta(G)+1\), and
\(\noderef{MissingColorsNonemptyOfMaximumDegree}\) applied at \(b_3\)
supplies such a \(\gamma\).  Since the old first color \(a\)
is missing on the tree, the closedness condition
\(\noderef{ColoredVertexSetClosed}\) forbids both \(a\) and \(\gamma\)
on the boundary.  They are distinct from \(b\): \(b\) is not missing anywhere
on the tree, whereas \(\gamma\) is missing at \(b_3\).  If
\(\gamma=a\), the outside swap below is the identity and all subsequent
claims remain valid; otherwise apply
\(\noderef{OutsideTreeColorSwapStable}\) to \(a,\gamma\).

Put
\(\sigma=\varphi/(G-T,a,\gamma)\).  The cited stability lemma says that
\(\sigma\) is proper, leaves exactly \(e_0\) uncolored, and fixes every
edge incident with \(T\), as well as every missing-color set on \(T\).
Thus the coloring-independent clauses of the \(k\)-triple, together with
those two coloring clauses, make \((G,e_0,\sigma)\) a \(k\)-triple.
The displayed edges and all earlier missing witnesses in
\(\noderef{TashkinovTree}\) are unchanged, so \(T\) remains Tashkinov; the same
pointwise equality on incident edges preserves closedness.

Every boundary edge \(f_i\) remains colored \(b\).  By the exact
first-contact clause of \(\noderef{TreeExitPath}\), before the first tree
vertex of an exit leg no edge is incident with \(T\), so the outside
swap replaces \(a\) by \(\gamma\) and fixes \(b\) all along that portion.
The missing status at the remote end changes by the same transposition.
Consequently the same three vertex-edge sequences, regarded in the
\((\gamma,b)\)-incidence graph, are exit paths with the same supports and
lengths.  Support-disjointness, endpoint distinctness, and the exact
boundary witnesses \(f_i\) are unchanged.  These data define the asserted
counterexample \(D\), with \(D.a=\gamma\), and give all the displayed
equalities and order relations.
\end{proof}
